This paper designed, instantiated and evaluated an explainable decision support framework for multi-item FMCG inventory planning. Its decision model sets order quantities and safety stocks jointly under ordering, holding, expected-shortage and quantity-discounted purchasing costs, a procurement budget, warehouse capacities and service targets, and it is convex within each price regime, which yields a lower bound that certifies the quality of any recommended policy. An ABC algorithm adapted to the problem returned feasible policies in every run and, at an equal evaluation budget, the lowest-cost policies on most of the synthetic instances studied (with the lowest median also at equal computing time, though significantly so only against adapted GA, not for the smallest portfolios, and tied with adapted GA and PSO on small discounted instances), while generic penalty-based GA, PSO and ABC rarely returned feasible policies when the resources were binding; the ablation traces the gain to inventory-specific constraint handling and a log-scale perturbation rather than to the bee-colony search as such. Replaying the policies on held-out demand showed that model feasibility does not guarantee realised service, mainly because the continuous-review model ignores daily-review undershoot (modelling it reduced the products missing their targets from 27 to 9 of 50), which the decision artifact reports next to cost drivers, binding resources, the certified gap and rule-based recommendations whose thresholds were fixed in advance. The framework is implemented in a working decision support system in which planners retain and record the final decision. The evidence is synthetic and technical; application to the thesis' distributor data and a user study of the explanations are the next steps.
